We can therefore extend @@M0@@ in the evident way to a homomorphism of @@M1@@-algebras: @@M2@@. We therefore have:
\sgaStableTarget{sga6.e05.eq.6-3-3}@@M3@@

\sgaStableTarget{sga6.e05.def.6-4}\textbf{Definition 6.4.} \uline{The homomorphism} @@M4@@ \uline{just} \uline{defined is called the Chern character}.

If @@M5@@ is of characteristic @@M6@@ (i.e. is a @@M7@@-algebra), the Chern character is an isomorphism: indeed, let @@M8@@; we can write @@M9@@, where @@M10@@ has no term of degree zero; we then define an inverse morphism of @@M11@@ by @@M12@@.

\sgaStableTarget{sga6.e05.prop.6-5}\textbf{Proposition 6.5.} \uline{Let} @@M13@@ \uline{be a} \uline{graded} @@M14@@-\uline{algebra}; \uline{on} @@M15@@ \uline{consider the filtration defined as follows} (\uline{degree filtration}):
\sgaStableTarget{sga6.e05.eq.6-5-1}@@M16@@
\uline{where} @@M17@@. \uline{We have the following congruences}:

\begin{enumerate}[label=\roman*)]
\item \uline{Let} @@M18@@, \uline{and} @@M19@@; \uline{put} @@M20@@; \uline{then}:
\sgaStableTarget{sga6.e05.eq.6-5-2}@@M21@@

\item \uline{Let} @@M22@@, @@M23@@; \uline{put}
@@M24@@, @@M25@@; \uline{then}:
\sgaStableTarget{sga6.e05.eq.6-5-3}@@M26@@
\uline{and consequently}, \uline{the degree filtration is a ring filtration}.
\end{enumerate}

i) The Chern character is by construction a morphism of filtered functors; from the congruence @@M27@@, we deduce
